\section{Python implementation: tabular Q-learning in a small GridWorld}
\label{sec:gridworld-q}

\begin{kaggleproject}{GridWorld --- learn a path by trial and error}
\textbf{Question.} Can a Q-table learn to move from a start cell to a goal while avoiding a negative terminal state using only rewards from repeated interaction?
\end{kaggleproject}

Reinforcement learning is different from the earlier chapters because there is no static CSV to download: the agent must interact with an environment. To keep the first RL example self-contained, the script defines a tiny $4\times4$ GridWorld directly in NumPy. This avoids a Kaggle competition, avoids the \texttt{kaggle\_environments} package, and lets the Q-learning update remain completely visible.

The agent begins with an empty Q-table, explores with an $\varepsilon$-greedy policy, updates one state--action value after each move, gradually reduces exploration, and finally evaluates the learned greedy policy over repeated episodes.

\textbf{Before running.} Predict the qualitative pattern you expect from the experiment before you look at the output or plot.

\begin{labmode}
Stage 5: specify the environment change and predicted consequence before training. Treat the script as a transparent simulator you are expected to modify.
\end{labmode}

\textbf{Part A --- define the environment, Q-table, and behavior policy.}

\lstinputlisting[style=mlpython,firstline=1,lastline=39]{all_experiments/lab19-gridworld_qlearning.py}

\textbf{Part B --- train with Q-learning and inspect the learning curve.}

\lstinputlisting[style=mlpython,firstline=40,lastline=72]{all_experiments/lab19-gridworld_qlearning.py}

\textbf{Part C --- evaluate the learned greedy policy separately.}

\lstinputlisting[style=mlpython,firstline=74,lastline=91]{all_experiments/lab19-gridworld_qlearning.py}


\begin{pitfall}
A tabular state representation grows rapidly with environment complexity. GridWorld is intentionally small so every idea can be inspected directly; larger reinforcement-learning problems usually need better state representations or function approximation.
\end{pitfall}

\begin{labdeliverable}
Submit the reward specification, one predicted behavioral consequence before training, a learning curve, the final greedy-policy success rate over repeated evaluation episodes, and one controlled modification to reward, trap location, exploration, or discounting with an explanation of how the learned behavior changed.
\end{labdeliverable}
